\scanpage{25}
(ii) $\longrightarrow$ (ii bis) triviaux. Enfin \deleted{\unsure{(i bis)}}\inserted{(ii bis)} $\Rightarrow$ (ii), car
si $M$ est un module de type fini sur $A$ tel que
$\operatorname{Hom}(M,M)$ soit monogène, il résulte de Nakayama \unsure{que}
$\operatorname{id}_M$ en est un générateur.

\uline{Cor. 3} Conditions équivalentes

(i) $\unsure{g|\eta|\mathfrak D}$ est bijectif

(ii) $H^n(\Omega)$ est dualisant

(iii) $A\to\operatorname{Hom}(\Omega,\Omega)$ est un isom

(iv) $\operatorname{Hom}(\Omega,\Omega)\simeq A$

\deleted{\unsure{et}} Alors \deleted{\unsure{aussi}} les idéaux pr. associés à $A$
sont tous de \deleted{\unsure{dim|codim}} corang $n$.

\uline{Cor. 4} Les cond. du cor. 3 sont vérifiées
sous les conditions suivantes : a) $A$ est
Cohen-Macaulay \qquad (b) $A$ est normal.

\deleted{\unsure{(a)}} (a) Résulte du th. 4, cor. 3.

% slanted
\margin{à expliciter en \unsure{lemme}}
Quant à (b), comme les conditions du cor. 1
sont vérifiées ($A$ est intègre) il suffit
de prouver que $\operatorname{Hom}(\Omega,\Omega)\unsure{>|\supset} A$ est en
fait $=A$. Or $\Omega$ est sans torsion
[\deleted{ses idéaux pr. associés} $0$ est le seul idéal pr. associé,
car $A$ intègre] donc $\operatorname{Hom}(\Omega,\Omega)\subset K$ corps
% archive sheet number: 24

